% TIETOTEKNIIKAN KANDIDAATINTUTKIELMA
\documentclass[utf8,bachelor]{gradu3}
\usepackage[bookmarksopen,bookmarksnumbered,linktocpage]{hyperref}
\usepackage{graphicx}
\addbibresource{Lahteet.bib}

% META

\begin{document}

\title{Reaaliaikaisesti renderöidyn vektorigrafiikan käyttö videopeleissä}
\translatedtitle{Realtime rendered vector graphics in videogames}

\avainsanat{avain1, avain2, avain3}
\keywords{avainsanat englanniksi}
\tiivistelma{
    Tiivistelmä on tyypillisesti 5-10 riviä pitkä esitys työn pääkohdista (tausta, tavoite, tulokset, johtopäätökset).
}
\abstract{
    Englanninkielinen versio tiivistelmästä.
}

\author{Ville Kankaanpää}
\contactinformation{\texttt{vivakank@student.jyu.fi}}
%\supervisor{}

% Otsikko ja metatiedot
\maketitle

% Sisällysluettelo
\mainmatter

% Teksti

\chapter{Johdanto}

Videopelejä pelataan nykyään monilla erikokoisilla ja -resoluutioisilla näytöillä. Esimerkiksi Bethesda Softworksin julkaisema Fallout Shelter -peli on saatavilla älypuhelimille, tietokoneille ja pelikonsoleille \parencite{RefWorks:doc:5bd6d887e4b0a1f99c62e6de}. Käyttäessä perinteistä bittikarttagrafiikkaa kehittäjien täytyy tehdä kompromisseja kohdentaessaan pelejään eri tarkkuuksisille näytöille. Vaihtoehtoja ovat kuvan venyttäminen eri kokoihin, mikä tuottaa epätarkkoja tai pikselöityneitä kuvia, tai monen eri kokoisen kuvan tekeminen, mikä lisää työmäärä ja lopullisen pelin tiedostokokoa. \parencite{RefWorks:doc:5bd8319de4b03ae5c9b276b8}

Vektorigrafiikka on ratkaisu näihin ongelmiin. Se sallii kuvan venyttämisen eri kokoiseksi ilman laadun kärsimistä. Vektorigrafiikkatiedosto, kuten SVG, on usein paljon pienempi kooltaan kuin bittikarttatiedosto, kuten PNG \parencites{RefWorks:doc:5bdc5224e4b05afcfde5b159}{RefWorks:doc:5bdc5292e4b05afcfde5b171}. Kuviosta \ref{vertailu} voi huomata, että jos kuvan haluaa esittää skaalattuna alimman rivin kokoiseksi, vasta 512x512 pikselin -versiossa on tarpeeksi resoluutiota, eikä se näytä epäselvältä, toisin kuin sitä pienemmät versiot. 512x512 pikselin -version tiedostokoko on kuitenkin yli kaksinkertainen suhteessa alkuperäisen SVG-kuvan tiedostokokoon, ja jos kehittäjä haluaa pelin tukevan kaikkia kuviossa mainittuja resoluutioita tulee yhteenlasketuksi tiedostokooksi 612 kilotavua, yhdeksän kertaa alkuperäisen SVG-kuvan tiedostokoko.

SVG-tiedosto tukee myös animointia, ja yhdessä tiedostossa voi olla monta animaatiota \parencite{RefWorks:doc:5bd74719e4b0e42e08f6333b}. Tällä tavalla esimerkiksi pelihahmolle saa samaan tiedostoon toimeton-, juoksu- ja hyppyanimaation. Vektorigrafiikan renderöinti on perinteisesti ollut hidasta, eikä sitä ole voitu suorittaa tarpeeksi nopeasti sulavan liikkeen aikaansaamiseksi. Nykyiset teknologiat ja laitteistot kuitenkin mahdollistavat vektorikuvien esittämisen huomattavasti aiempaa nopeammin. Vektorikuvia on myös mahdollista esittää 3D-gafiikan seassa. \parencite{RefWorks:doc:5bc4a5cce4b080e02f7eff1b} Kysymykseksi jääkin, miksei reaaliajassa renderöidyn vektorigrafiikan käyttö videopeleissä ole yleistynyt, vaikka teknologia sen sallii.

Tässä kandidaatintutkielmassa tutkin vektorigrafiikoiden käyttöä videopeleissä kirjallisuuskatsauksen metodein. Pohdin vektorigrafiikoiden etuja ja haittoja verrattuna bittikarttagrafiikoihin 2D-pelien kontekstissa, mutta myös 3D-objektien teksturointia sivuten. Teen myös katsauksen peleissä käytetyn vektorigrafiikan historiasta, nykytilasta ja mahdollisista tulevaisuuden näkymistä.

\begin{figure}[!b]
	\centering
	\includegraphics[width=1.0\textwidth]{Ghostscript_Tiger_Comparison.pdf}
    \caption{SVG-kuvan ja erikokoisten PNG-kuvien vertailua. Ylemmällä rivillä on alkuperäinen kuva ja alemmalla suurennos yksityiskohdasta.}\label{vertailu}
\end{figure}

\chapter{Vektorigrafiikka}

\section{Vektorigrafiikat ja niiden käyttö}

% - Määrittely
% - Ominaisuudet
% - Käyttökohteet
%     * Fontit ja logot
%     * Datan visualisointi ja pohjapiirustukset
%     * Käyttöliittymät (?)
%     * Pelien assetit
%         ~ Banner Saga
%         ~ Guacamelee
%         ~ Scribblenauts

Vektorigrafiikat ovat tietynlaisia tietokoneella generoituja kuvia. Toisin kuin bittikarttagrafiikat, jotka sisältävät jokaisen pikselin väriarvot, vektorigrafiikat sisältävät tietoa pisteistä ja niiden väliin tulevista suorista viivoista, käyristä viivoista ja muista muodoista \parencites{RefWorks:doc:5bdc5224e4b05afcfde5b159}{RefWorks:doc:5bdc5292e4b05afcfde5b171}. Erilaista tietosisältöä ovat esimerkiksi viivan väri, paksuus, katkonaisuus ja käyttäytyminen kulmissa (terävät, pyöreät tai tylpät kulmat). Viivan sisäpuolen voi täyttää yhtenäisellä värillä, liukuvärillä tai kuviolla, ja sitä voi myös leikata toisella viivalla. \parencite{RefWorks:doc:5bd74719e4b0e42e08f6333b} Koska vektorigrafiikan pisteitä ei ole sidottu pikseleihin vaan omaan koordinaatistoonsa, voidaan sitä skaalata minkä kokoiseksi tahansa ilman että laatu kärsii tai tiedostokoko muuttuu \parencite{RefWorks:doc:5bdc5292e4b05afcfde5b171}.

Käytettyjä vektorigrafiikan tiedostomuotoja ovat esimerkiksi Encapsulated PostScript (EPS), Flash (SWF), Adobe Illustrator (AI) ja Scalable Vector Graphics (SVG) \parencite{RefWorks:doc:5bdc5224e4b05afcfde5b159}. Tässä tutkielmassa keskityn vektorigrafiikan tarkasteluun SVG-tiedostomuodon pohjalta, sillä se on W3C:n standardi vektorigrafiikoille. SVG:n spesifikaatio on myös kattavampi kuin esimerkiksi EPS:n \parencite{RefWorks:doc:5bdc900de4b0afdabde32fb7} ja SWF:n \parencite{RefWorks:doc:5bdc9306e4b069b454d8459e}. SVG-tiedostomuoto tukee edellä mainittujen ominaisuuksien lisäksi muun muassa bittikarttakuvien upottamista, tekstien lisäämistä ja animointia. \parencite{RefWorks:doc:5bd74719e4b0e42e08f6333b}

Vektorigrafiikoilla on useita käyttökohteita. \textcite{RefWorks:doc:5bdc5224e4b05afcfde5b159} esittelee käyttökohteiksi muun muassa logot, tekniset piirustukset ja kaaviokuvat. \textcite{RefWorks:doc:5bdd9b5ae4b0954dddb26329} lisäävät käyttökohteisiin esimerkiksi fontit, kartat, käyttöliittymät ja pelit. Seuraavissa kappaleissa esittelen perusteluita miksi vektorigrafiikat soveltuvat hyvin näihin käyttötarkoituksiin.

Yhdestä vektorikuvasta on helppo luoda erikokoisia versioita \parencite{RefWorks:doc:5bdc5224e4b05afcfde5b159}. Esimerkiksi yritysten logot voivat esiintyä monessa eri koossa pienistä käyntikorteista suuriin tienvarsimainoksiin. On siis tärkeää, että logosta on olemassa useita eri resoluutioisia versioita, tai yksi resoluutioriippumaton versio, vektorikuva. Samasta syystä kuin logojen, teksteissä käytettyjen fonttien tulee olla hyvin skaalautuvia. Nykyisin suurimmassa osassa fontteja kirjaimia esittävät kuvat ovatkin juuri vektorikuvia, mikä sallii tekstin esittämisen eri fonttikoilla \parencite{RefWorks:doc:5bdc5224e4b05afcfde5b159}. Vektorikuvaa on myös helpompi muokata kuin bittikarttakuvaa, sillä sen elementit ovat itsenäisiä objekteja \parencite{RefWorks:doc:5bdc5224e4b05afcfde5b159}. Jos yritys haluaa vaihtaa esimerkiksi värejä, fontteja tai elementtien sijoittelua logossaan, onnistuu se vektorikuvan kanssa helposti.

Vektorigrafiikoilla voidaan esittää laadukkaita interaktiivisia kuvaajia \parencite{RefWorks:doc:5bdc736ae4b05afcfde5d6e4}. Artikkelissaan Lang ja Nolan esittelevät R-ohjelmointikielellä tekemäänsä työkalua, jolla pystyy luomaan datasta erinäköisiä SVG-kuvia. Kuvaan pystyy SVG-tiedoston ominaisuuksien johdosta lisäämään interaktiivisia toiminnallisuuksia. Esimerkiksi elementtejä saa animoitua, kun niitä osoitetaan tai klikataan hiirellä, ja liukusäätimestä voidaan säätää kuvassa esitettävää sisältöä. \parencite{RefWorks:doc:5bdc736ae4b05afcfde5d6e4} Myös Googlen Charts API muodostaa SVG-muotoisia interaktiivisia kaavioita sille syötetystä datasta \parencite{RefWorks:doc:5bde1508e4b0afdabde34e20}.

Tietokoneen käyttöliittymät voivat olla toteutettu vektorigrafiikkana. Microsoftin Windows Presentation Foundationia (WPF) käyttävät ohjelmat ovat vektoripohjaisia. WPF lupaa, että kehittäjän ei tarvitse murehtia päätelaitteen monitorin koosta tai resoluutiosta, sillä se hoitaa skaalaamisen sopivan kokoiseksi. \parencite{RefWorks:doc:5bde3d66e4b0a1f99c63ae5a} Unixin työpöytäympäristö KDE 4:n grafiikat ovat vektorigrafiikoita. Työpöytä ja siihen liittyvät elementit ovat joukkoja SVG-kuvia ja niitä muokkaamalla voidaan säätää värejä ja muotoja. Kaikki ikonit ovat SVG-kuvia, mutta reaaliaikaisen renderöinnin sijasta ikonista pyydetään tietynkokoista versiota, joka rasteroidaan välimuistiin. Tällä tavalla parannetaan suorituskykyä, jota vektorikuvan jatkuva uudelleenrenderöinti heikentäisi. \parencite{RefWorks:doc:5bded7f3e4b09dfad48b846c}

Vektorigrafiikoita käytetään peleissä eri tavoin. \textcite{RefWorks:doc:5bdc818be4b0965708b5a0f0} kertoo artikkelissaan kuinka Guacamelee!-pelin eri vaiheissa käytettiin vektorigrafiikoita. Ensimmäiset pelin luonnokset olivat Flashilla tehtyjä vektorianimaatioita. Myös pelin lopulliset grafiikat tehtiin Flashilla, mikä takasi sulavat animaatiot ja säästi muistia. \parencite{RefWorks:doc:5bdc818be4b0965708b5a0f0} \textcite{RefWorks:doc:5bdc7473e4b05afcfde5d742} tekemässä haastattelussa Brittany Aubert kertoo Scribblenauts Unlimited -pelin grafiikoiden olevan vektorikuvia, jotta objekteja voisi venyttää ilman kuvanlaadun kärsimistä. \textcite{RefWorks:doc:5bd88176e4b0afdabde2b8e2} puolestaan kertoo artikkelissan Banner Saga -pelin animointiprosessista. Animoitavat liikkeet kuvattiin ihmisten näyttelemänä ja rotoskoopattiin Flashilla sarjaksi vektorikuvia. Sen jälkeen Flash-animaatiot rasteroitiin joukoksi bittikarttakuvia. \parencite{RefWorks:doc:5bd88176e4b0afdabde2b8e2}

\section{Vektorigrafiikan ja bittikarttagrafiikan eduista ja haitoista}

% - Level of Detail
% - Renderöinti
% - Tiedostokoko
% - Muokkaus
% - Animointi

\chapter{Vektorigrafiikat videopeleissä}

\section{Vektorigrafiikan historia peleissä}

% - Oikeat vektoripelit
% - Another World
% - Flash-selainpelit

\section{Vektorigrafiikan nykytila peleissä}

% - Banner Saga
% - Guacamelee
% - Scribblenauts
% - TF2

\section{Ongelmia vektorigrafiikan käytössä}

% - Ei virallista engine tukea
% - "Rajoittunut" tyyli

\chapter{Pohdintaa vektorigrafiikoiden tulevaisuudesta}

\section{Selainpelit}

% - Selaimessa toimivien sovellusten kasvu
% - Suora tuki SVG:lle ja animaatioille

\section{Vektorigrafiikka 3D-kappaleiden tekstuureissa}

% - Supertekstuurit
% - TF2

\chapter{Yhteenveto}

% - Tarvitaan esimerkki proto

\printbibliography

\end{document}